\documentclass[10pt,a4paper]{amsart}
\usepackage{amsmath,amssymb,amsthm,graphicx}
% Retained preamble for the Stacks Project chapter snapshot in this directory.
% The chapter's own preamble.tex is a copy of the shared Stacks Project preamble; the block
% below is sliced MECHANICALLY, line for line, from it. Nothing was added, removed or
% redefined; the wrapper class, the xr-hyper/hyperref wiring and the externaldocument list
% are omitted because the assembler supplies the document class.
%
% lines taken from preamble.tex: 3, 5, 9, 11, 12, 156, 157, 158, ... (51 lines total)
%
% package coverage restored from the chapter's own preamble (2026-10-13): xy, xy-2cell

\usepackage[all]{xy}
\xyoption{2cell}
\UseAllTwocells

\newenvironment{reference}{\comment}{\endcomment}
\newenvironment{slogan}{\comment}{\endcomment}
\newenvironment{history}{\comment}{\endcomment}
\theoremstyle{plain}
\newtheorem{theorem}[subsection]{Theorem}
\newtheorem{proposition}[subsection]{Proposition}
\newtheorem{lemma}[subsection]{Lemma}
\theoremstyle{definition}
\newtheorem{definition}[subsection]{Definition}
\newtheorem{example}[subsection]{Example}
\newtheorem{exercise}[subsection]{Exercise}
\newtheorem{situation}[subsection]{Situation}
\theoremstyle{remark}
\newtheorem{remark}[subsection]{Remark}
\newtheorem{remarks}[subsection]{Remarks}
\numberwithin{equation}{subsection}
\def\lim{\mathop{\mathrm{lim}}\nolimits}
\def\colim{\mathop{\mathrm{colim}}\nolimits}
\def\Spec{\mathop{\mathrm{Spec}}}
\def\Hom{\mathop{\mathrm{Hom}}\nolimits}
\def\Ext{\mathop{\mathrm{Ext}}\nolimits}
\def\SheafHom{\mathop{\mathcal{H}\!\mathit{om}}\nolimits}
\def\SheafExt{\mathop{\mathcal{E}\!\mathit{xt}}\nolimits}
\def\Sch{\mathit{Sch}}
\def\Mor{\mathop{\mathrm{Mor}}\nolimits}
\def\Ob{\mathop{\mathrm{Ob}}\nolimits}
\def\Sh{\mathop{\mathit{Sh}}\nolimits}
\def\NL{\mathop{N\!L}\nolimits}
\def\CH{\mathop{\mathrm{CH}}\nolimits}
\def\proetale{{pro\text{-}\acute{e}tale}}
\def\etale{{\acute{e}tale}}
\def\QCoh{\mathit{QCoh}}
\def\Ker{\mathop{\mathrm{Ker}}}
\def\Im{\mathop{\mathrm{Im}}}
\def\Coker{\mathop{\mathrm{Coker}}}
\def\Coim{\mathop{\mathrm{Coim}}}
\DeclareMathSymbol{\boxtimes}{\mathbin}{AMSa}{"02}
\def\QCohstack{\mathcal{QC}\!\mathit{oh}}
\def\Cohstack{\mathcal{C}\!\mathit{oh}}
\def\Spacesstack{\mathcal{S}\!\mathit{paces}}
\def\Quotfunctor{\mathrm{Quot}}
\def\Hilbfunctor{\mathrm{Hilb}}
\def\Curvesstack{\mathcal{C}\!\mathit{urves}}
\def\Polarizedstack{\mathcal{P}\!\mathit{olarized}}
\def\Complexesstack{\mathcal{C}\!\mathit{omplexes}}
\def\Pic{\mathop{\mathrm{Pic}}\nolimits}
\def\Picardstack{\mathcal{P}\!\mathit{ic}}
\def\Picardfunctor{\mathrm{Pic}}
\def\Deformationcategory{\mathcal{D}\!\mathit{ef}}
\usepackage{cleveref}
\crefname{originallemma}{original lemma}{original lemmas}
\Crefname{originallemma}{Original lemma}{Original lemmas}
\makeatletter
\newlabel{bootstrap-theorem-final-bootstrap}{{04S6}{source}{}{sourcepointer-1}{}}
\newlabel{bootstrap-theorem-final-bootstrap@cref}{{[originallemma][0][]04S6}{[1][0][]source}{}{}{}}
\makeatother
\crefname{originallemma}{original lemma}{original lemmas}
\Crefname{originallemma}{Original lemma}{Original lemmas}
\makeatletter
\newlabel{equation-beauville-laszlo-glueing-spaces}{{0F9N}{source}{}{sourcepointer-2}{}}
\newlabel{equation-beauville-laszlo-glueing-spaces@cref}{{[originallemma][0][]0F9N}{[1][0][]source}{}{}{}}
\makeatother
\crefname{originallemma}{original lemma}{original lemmas}
\Crefname{originallemma}{Original lemma}{Original lemmas}
\makeatletter
\newlabel{lemma-glueing-affines}{{0F9Q}{source}{}{sourcepointer-3}{}}
\newlabel{lemma-glueing-affines@cref}{{[originallemma][0][]0F9Q}{[1][0][]source}{}{}{}}
\makeatother
\crefname{originallemma}{original lemma}{original lemmas}
\Crefname{originallemma}{Original lemma}{Original lemmas}
\makeatletter
\newlabel{lemma-glueing-affines-etale}{{0F9R}{source}{}{sourcepointer-4}{}}
\newlabel{lemma-glueing-affines-etale@cref}{{[originallemma][0][]0F9R}{[1][0][]source}{}{}{}}
\makeatother
\crefname{originallemma}{original lemma}{original lemmas}
\Crefname{originallemma}{Original lemma}{Original lemmas}
\makeatletter
\newlabel{properties-lemma-quasi-affine-presentation}{{0F82}{source}{}{sourcepointer-5}{}}
\newlabel{properties-lemma-quasi-affine-presentation@cref}{{[originallemma][0][]0F82}{[1][0][]source}{}{}{}}
\makeatother
\crefname{originallemma}{original lemma}{original lemmas}
\Crefname{originallemma}{Original lemma}{Original lemmas}
\makeatletter
\newlabel{spaces-descent-lemma-descending-property-closed-immersion}{{0420}{source}{}{sourcepointer-6}{}}
\newlabel{spaces-descent-lemma-descending-property-closed-immersion@cref}{{[originallemma][0][]0420}{[1][0][]source}{}{}{}}
\makeatother
\title{Beauville-Laszlo glueing realizes every quasi-affine triple by a separated space}
\author{The Stacks Project Authors}
\date{stacks-project:spaces-pushouts.tex:local-snapshot-20261012; GNU FDL 1.2 or later}
\begin{document}
\maketitle
\noindent\textit{Editorial note.} Statement span is lines 2757-2764 of sources/stacks-spaces-pushouts/pinned-original-source.tex; proof span is lines 2766-2844, the complete author proof (77 body lines).\par\medskip
\noindent\textit{Editorial note.} Navigation stubs: 6 cited-but-not-included labels are supplied as non-load-bearing pointers with their pinned tags where the snapshot tag index has them.\par\medskip
\noindent\textit{Editorial note.} extra\_packages is empty on purpose: the retained preamble already loads xy with options and a second declaration clashes in the assembler wrapper.\par\medskip
\noindent\textit{Editorial note.} Text is retained verbatim, including the author's own typos and placeholders.\par\medskip
\noindent\textit{Editorial note.} Queue index 7 of queue\_spaces-pushouts.json (pb=2899, stmt 2757-2764, proof 2766-2844).\par\medskip
\noindent\textit{Editorial note.} spaces-pushouts wave-1: fresh adjudication accept (2026-10-14); six-check dispatched pool (§2 row appended 02:4x).\par\medskip
\noindent\textit{Allowed standard prerequisites (editorial).}\begin{itemize}
\item Explicit prerequisite boundary. The cited inputs are prerequisites; the argument in the retained span is not replaced by any cited result.
\end{itemize}\medskip
\ifdefined\hypertarget\hypertarget{sourcepointer-1}{}\fi
\noindent\textit{Source reference.} Original lemma 04S6: sources/stacks-spaces-pushouts/pinned-original-source.tex, source lines 2844--2844. Cited by the retained statement or proof but not part of the collected spans; supplied only as a navigation pointer. The argument does not depend on its content.\par\medskip
\ifdefined\hypertarget\hypertarget{sourcepointer-2}{}\fi
\noindent\textit{Source reference.} Original lemma 0F9N: sources/stacks-spaces-pushouts/pinned-original-source.tex, source lines 2470--2474. Cited by the retained statement or proof but not part of the collected spans; supplied only as a navigation pointer. The argument does not depend on its content.\par\medskip
\ifdefined\hypertarget\hypertarget{sourcepointer-3}{}\fi
\noindent\textit{Source reference.} Original lemma 0F9Q: sources/stacks-spaces-pushouts/pinned-original-source.tex, source lines 2547--2555. Cited by the retained statement or proof but not part of the collected spans; supplied only as a navigation pointer. The argument does not depend on its content.\par\medskip
\ifdefined\hypertarget\hypertarget{sourcepointer-4}{}\fi
\noindent\textit{Source reference.} Original lemma 0F9R: sources/stacks-spaces-pushouts/pinned-original-source.tex, source lines 2582--2589. Cited by the retained statement or proof but not part of the collected spans; supplied only as a navigation pointer. The argument does not depend on its content.\par\medskip
\ifdefined\hypertarget\hypertarget{sourcepointer-5}{}\fi
\noindent\textit{Source reference.} Original lemma 0F82: sources/stacks-spaces-pushouts/pinned-original-source.tex, source lines 2844--2844. Cited by the retained statement or proof but not part of the collected spans; supplied only as a navigation pointer. The argument does not depend on its content.\par\medskip
\ifdefined\hypertarget\hypertarget{sourcepointer-6}{}\fi
\noindent\textit{Source reference.} Original lemma 0420: sources/stacks-spaces-pushouts/pinned-original-source.tex, source lines 2844--2844. Cited by the retained statement or proof but not part of the collected spans; supplied only as a navigation pointer. The argument does not depend on its content.\par\medskip
\begin{lemma}
\label{lemma-glueing-quasi-affines}
Let $(R \to R', f)$ be a glueing pair, see above. Any object
$(V, V', Y')$ of $\textit{Spaces}(U \leftarrow U' \to X')$
with $V$, $V'$, $Y'$ quasi-affine is isomorphic to
the image under the functor (\ref{equation-beauville-laszlo-glueing-spaces})
of a separated algebraic space $Y$ over $X$.
\end{lemma}

\begin{proof}
Choose $n'$, $T' \to Y'$ and $n_1$, $T_1 \to V$ as in
Properties, Lemma \ref{properties-lemma-quasi-affine-presentation}.
Picture
$$
\xymatrix{
& &
T_1 \times_V V' \times_Y T' \ar[ld] \ar[rd] \\
T_1 \ar[d] &
T_1 \times_V V' \ar[l] \ar[dr] & &
V' \times_{Y'} T' \ar[r] \ar[dl] &
T' \ar[d] \\
V & &
V' \ar[rr] \ar[ll] & &
Y'
}
$$
Observe that $T_1 \times_V V'$ and $V' \times_{Y'} T'$
are affine (namely the morphisms $V' \to V$ and $V' \to Y'$
are affine as base changes of the affine morphisms $U' \to U$
and $U' \to X'$). By construction we see that
$$
\mathbf{A}^{n'}_{T_1 \times_V V'} \cong
T_1 \times_V V' \times_{Y'} T' \cong
\mathbf{A}^{n_1}_{V' \times_{Y'} T'}
$$
In other words, the affine schemes $\mathbf{A}^{n'}_{T_1}$
and $\mathbf{A}^{n_1}_{T'}$ are part of a triple making an affine object of
$\textit{Spaces}(U \leftarrow U' \to X')$.
By Lemma \ref{lemma-glueing-affines}
there exists a morphism of affine schemes $T \to X$
and isomorphisms $U \times_X T \cong \mathbf{A}^{n'}_{T_1}$
and $X' \times_X T \cong \mathbf{A}^{n_1}_{T'}$ compatible
with the isomorphisms displayed above.
These isomorphisms produce morphisms
$$
U \times_X T \longrightarrow V
\quad\text{and}\quad
X' \times_X T \longrightarrow Y'
$$
satisfying the property of
Properties, Lemma \ref{properties-lemma-quasi-affine-presentation}
with $n = n' + n_1$ and moreover define a morphism from the triple
$(U \times_X T, U' \times_X T, X' \times_X T)$ to
our triple $(V, V', Y')$ in the category
$\textit{Spaces}(U \leftarrow U' \to X')$.

\medskip\noindent
By Lemma \ref{lemma-glueing-affines} there is an affine scheme $W$ whose
image in $\textit{Spaces}(U \leftarrow U' \to X')$ is isomorphic to
the triple
$$
((U \times_X T) \times_V (U \times_X T),
(U' \times_X T) \times_{V'} (U' \times_X T),
(X' \times_X T) \times_{Y'} (X' \times_X T))
$$
By fully faithfulness of this construction, we obtain
two maps $p_0, p_1 : W \to T$ whose base changes
to $U, U', X'$ are the projection morphisms.
By Lemma \ref{lemma-glueing-affines-etale}
the morphisms $p_0, p_1$ are flat and of finite presentation and
the morphism $(p_0, p_1) : W \to T \times_X T$ is a closed immersion.
In fact, $W \to T \times_X T$ is an equivalence relation: by the lemmas
used above we may check symmetry, reflexivity, and transitivity
after base change to $U$ and $X'$, where these are obvious (details omitted).
Thus the quotient sheaf
$$
Y = T/W
$$
is an algebraic space for example by
Bootstrap, Theorem \ref{bootstrap-theorem-final-bootstrap}.
Since it is clear that $Y/X$ is sent to the triple $(V, V', Y')$.
The base change of the diagonal $\Delta : Y \to Y \times_X Y$
by the quasi-compact surjective flat morphism $T \times_X T \to Y \times_X Y$
is the closed immersion $W \to T \times_X T$. Thus $\Delta$
is a closed immersion by Descent on Spaces, Lemma
\ref{spaces-descent-lemma-descending-property-closed-immersion}.
Thus the algebraic space $Y$ is separated and the proof is complete.
\end{proof}

\end{document}
